\documentclass{article}
\usepackage[UTF8]{ctex}
\usepackage{amsmath,amssymb,amsfonts}
\usepackage{enumitem}
\usepackage{geometry}
\geometry{a4paper, margin=1in}
% === 错题本排版优化（不参与题目标识） ===
\setlength{\parindent}{0pt}
\setlength{\parskip}{0.35em plus 0.10em minus 0.08em}
\setlength{\emergencystretch}{3em}
\setlist{itemsep=0.25em, topsep=0.40em, parsep=0pt}
\allowdisplaybreaks[2]
% === 错题本排版优化结束 ===
\begin{document}
\section*{华南师范大学2025年数学分析真题}
\section*{160．华南师范大学2025年数学分析试卷}

\subsection{题目1}
1． \(\lim _{n \rightarrow \infty} \frac{\sqrt{n+5}-\sqrt{n}}{\sqrt{n}-\sqrt{n+2005}}\) ．

\subsection{题目2}
2． \(\lim _{n \rightarrow \infty}\left(\frac{n+2}{n-1}\right)^{n}\) ．

\subsection{题目3}
3． \(\lim _{x \rightarrow+\infty}(\sin \sqrt{x+1}-\sin \sqrt{x})\) ．

\subsection{题目4}
4． \(\lim _{x \rightarrow 0} \frac{x \cos x-\sin x}{x^{2} \sin x}\) ．

\subsection{题目5}
5． \(\lim _{(x, y) \rightarrow(0,0)} \frac{x y}{\sqrt{x y+1}-1}\) ．

\subsection{题目6}
6．设 \(n\) 为正整数，\(y=\frac{x^{n}}{2-x}\) ，求 \(y^{(n)}\) ．

\subsection{题目7}
7．设 \(u=f\left(\frac{y}{x}+\frac{z}{y}+\frac{x}{z}\right)\) ，其中 \(f\) 二阶可导，求 \(\frac{\partial u}{\partial x}, \frac{\partial^{2} u}{\partial x \partial y}\) ．

\subsection{题目8}
8 ．求 \(\int e^{-x}\left(\frac{1+x}{1+x^{2}}\right)^{2} \mathrm{~d} x\) ．

\subsection{题目9}
9．计算 \(\iint_{D}|\sin (x+y)| \mathrm{d} x \mathrm{~d} y, D=[0, \pi] \times[0, \pi]\) ．

\subsection{题目10}
10．对于二元函数 \(f(x, y)\) ，如果两个累次极限存在，且

\[
\lim _{x \rightarrow 0} \lim _{y \rightarrow 0} f(x, y)=\lim _{y \rightarrow 0} \lim _{x \rightarrow 0} f(x, y) .
\]

问 \(\lim _{(x, y) \rightarrow(0,0)} f(x, y)\) 是否一定存在？请说明理由．

\subsection{题目11}
11．（10 分）设 \(f_{n}(x)=\frac{x^{2}}{1+n^{3} x^{3}}, n=1,2, \cdots\) ，试证明函数列 \(\left\{f_{n}(x)\right\}\) 在 \((0,+\infty)\) 内一致收敛．

\subsection{题目12}
12．（20 分）判断下列级数或反常积分的敛散性．

\begin{itemize}
  \item[（1）] \(\sum_{n=1}^{\infty} \frac{1}{n} \sin \frac{1}{n^{q}}, q \geq 0\) ．
  \item[（2）] \(\int_{3}^{+\infty} \frac{\ln x \mathrm{~d} x}{(x-\sqrt{x}+2)^{q}}, q>0\) ．
\end{itemize}

\subsection{题目13}
13．（10 分）求函数 \(f(x)=\frac{4}{x^{2}-2 x-3}\) 在 \(x=2\) 处的幂级数展开式．

\subsection{题目14}
14．（10 分）设在矩形区域 \(D=[a, b] \times[c, d]\) 上，\(F(x, y)\) 关于 \(x\) 连续，\(F_{y}(x, y)\) 存在且有界．证明： \(F(x, y)\) 在 \(D\) 上连续．
\end{document}
